\documentclass[conference,a4paper]{IEEEtran}
\IEEEoverridecommandlockouts
% The preceding line is only needed to identify funding in the first footnote. If that is unneeded, please comment it out.
\usepackage{cite}
\usepackage{amsmath,amssymb,amsfonts}
\usepackage{algorithmic}
\usepackage{graphicx}
\usepackage{textcomp}
\usepackage{xcolor}
\def\BibTeX{{\rm B\kern-.05em{\sc i\kern-.025em b}\kern-.08em
    T\kern-.1667em\lower.7ex\hbox{E}\kern-.125emX}}
% Project additions to the IEEE template.
\usepackage[T1]{fontenc}
\usepackage[utf8]{inputenc}
\usepackage{booktabs}
\usepackage{url}
\usepackage{etoolbox}
\patchcmd{\abstract}{---}{: }{}{\errmessage{Abstract punctuation patch failed}}
\patchcmd{\IEEEkeywords}{---}{: }{}{\errmessage{Keywords punctuation patch failed}}
\newcommand{\NaR}{\mathrm{NaR}}
\newcommand{\code}[1]{\texttt{#1}}
\usepackage[hidelinks]{hyperref}
\begin{document}

\title{An Exact Functional Model of Linear Takum Arithmetic\\for a RISC-V Vector Extension: Phase I}
\author{\IEEEauthorblockN{Arthur dos Santos Fernandes, Bruno Lima Soares,\\
Lucas Rodrigues Prosperi Ayala, Pedro Henrique Campos Lima,\\
and Pedro Lopes Ozelin}
\IEEEauthorblockA{Department of Electronic Engineering\\
Universidade Federal de Minas Gerais, Belo Horizonte, Brazil\\
Architecture and Computer Organization: 2026/2}}

\maketitle

\begin{abstract}
This report presents the first milestone of a project on Takum arithmetic for a
simplified RISC-V vector extension. We implement a scalar functional model of
linear Takum at 8, 16, and 32 bits, including field extraction, encoding,
decoding, width conversion, addition, subtraction, and multiplication. Exact
rational intermediates separate arithmetic correctness from host floating-point
precision. The rounding policy is explicitly defined as nearest/even on the
encoded bit stream, with saturation for finite values outside the representable
range. Validation combines exhaustive T8 arithmetic against an independent
boundary-table oracle, exhaustive T8/T16 round trips, and differential testing
against a pinned C reference implementation with a documented portability
correction. The local runs pass 20 tests and 548,793 differential comparisons.
This milestone establishes an arithmetic foundation; vector execution and
comparative sparse linear-algebra experiments remain future work.
\end{abstract}
\begin{IEEEkeywords}
Takum, tapered precision, computer arithmetic, functional simulation, RISC-V
\end{IEEEkeywords}

\section{Introduction}
Low-precision arithmetic requires explicit choices about dynamic range,
precision, exceptional values, and rounding. These choices affect both numerical
results and the architecture exposed to software. The AOC project proposes
studying a uniform family of Takum operands within a simplified vector model.
Its first milestone comprises the numerical representation and basic arithmetic;
the subsequent milestones introduce vector instructions and experimental
comparison with other number formats.

We select the \emph{linear} Takum variant. Hunhold's original work distinguishes
linear and logarithmic interpretations of a shared field structure~\cite{takum}.
The SIMD case study explicitly considers linear Takum in its dynamic-range
comparison~\cite{simd}. Because the course guidelines do not select a variant,
we state this choice throughout the implementation and documentation.

The contribution of this milestone is an inspectable arithmetic module that
implements the fields directly, together with reproducible evidence for its
behavior. The reference C library is used only in external tests. This report
makes no claim about processor speedup, silicon area, power consumption, or
numerical superiority over IEEE formats.

\section{Related Work}
Hunhold defines the Takum encoding and analyzes its properties, including
monotonicity and two's-complement negation. Definition 8 specifies the linear
variant; Algorithms 4 and 5 connect its fields to a conventional binary
floating-point representation~\cite{takum}. Those definitions provide the
mathematical basis for our codec.

The later SIMD study motivates a consistent operand family across precisions
and evaluates possible changes to AVX10.2~\cite{simd}. Its architectural
conclusions motivate investigation in other instruction sets; they do not
establish corresponding RISC-V performance results for this project.

The public \code{libtakum} C99 library provides a separately implemented
reference~\cite{libtakum}. Version 2.0.0 uses \code{takum8}, \code{takum16}, and
\code{takum32} for linear types. We pin a specific revision and preserve a
portability patch and both the original and corrected validation results.

\section{Representation and Numerical Contract}
\subsection{Field decoding}
An encoding contains a sign bit $S$, a direction bit $D$, three regime bits
interpreted as $R$, $r$ characteristic bits interpreted as $C$, and $p$ fraction
bits interpreted as $F$. The characteristic width is
\begin{equation}
 r=\begin{cases}R,&D=1,\\7-R,&D=0.\end{cases}
\end{equation}
The characteristic and exponent are
\begin{align}
 c&=\begin{cases}2^r-1+C,&D=1,\\1-2^{r+1}+C,&D=0,\end{cases}\\
 e&=(1-2S)(c+S).
\end{align}
When all characteristic bits are stored, $p=n-5-r$ and $f=F2^{-p}$. A finite,
nonzero value is
\begin{equation}
 x=\bigl[(1-3S)+f\bigr]2^e.
\end{equation}
For short encodings that end inside $C$, missing characteristic bits are
interpreted as trailing zeros, or \emph{ghost bits}. The software reports the
number of inserted zeros. There is no stored fraction in these cases.

All-zero bits encode zero. The sign bit followed by zeros encodes $\NaR$.
Table~\ref{tab:examples} gives selected T8 examples. The full 256-entry table is
exported as a CSV file. Negation is two's complement modulo $2^n$; zero and
$\NaR$ remain unchanged.

\begin{table}[t]
\caption{Selected linear T8 representations}
\label{tab:examples}
\centering
\begin{tabular}{lll}
\toprule
Hex & Binary & Exact value\\
\midrule
\code{00} & \code{00000000} & $0$\\
\code{01} & \code{00000001} & $2^{-239}$\\
\code{40} & \code{01000000} & $1$\\
\code{44} & \code{01000100} & $3/2$\\
\code{48} & \code{01001000} & $2$\\
\code{7f} & \code{01111111} & $2^{239}$\\
\code{80} & \code{10000000} & $\NaR$\\
\code{c0} & \code{11000000} & $-1$\\
\code{ff} & \code{11111111} & $-2^{-239}$\\
\bottomrule
\end{tabular}
\end{table}

\subsection{Encoding and rounding}
Finite inputs are represented as exact fractions. The encoder computes the
binary exponent from integer bit lengths and an exact comparison, constructs
$D$, $R$, and $C$, and quantizes the remaining bit stream. Positive magnitude is
encoded first; negative values then use two's-complement negation.

We adopt nearest/even rounding of the bit stream, following the implementation
policy in \code{libtakum} 2.0.0. This is an explicit project decision: the original
paper's Algorithm 1 increments based on the first discarded bit and does not
specify this tie-to-even refinement. Tests distinguish exact ties from values
arbitrarily close to them using rational remainders.

When truncation reaches the characteristic, bit-stream rounding need not select
the closest real value. For example, adjacent T8 codes \code{60} and \code{61}
represent $2^{15}$ and $2^{17}$. Their encoding threshold is $2^{16}$, whereas
the arithmetic midpoint of their values is $81920$. This difference is part of
the selected numerical contract and is tested explicitly.

Finite nonzero values below the minimum magnitude saturate to the minimum
nonzero magnitude, preserving sign. Values above the maximum magnitude saturate
to the corresponding finite maximum. Thus finite saturation produces neither
zero nor $\NaR$. The exact value zero remains zero. Nonfinite external inputs map
to $\NaR$, which propagates through all three arithmetic operations, including
multiplication by zero.

\section{Methodology and Proposed Architecture}
\subsection{Functional organization}
The implementation uses Python 3.10 or newer and only the standard library.
The arithmetic path is illustrated in Fig.~\ref{fig:architecture}. The codec
and ALU expose raw unsigned bit patterns at their interfaces; exact fractions
are internal values. The CLI additionally accepts decimal strings or fractions,
which avoids an unintended conversion through binary64.

\begin{figure}[t]
\centering
\fbox{\begin{minipage}{0.85\columnwidth}
\centering
T8 / T16 / T32 operand bits\\[3pt]
$\downarrow$\\[3pt]
Decode S, D, R, C, F and special values\\[3pt]
$\downarrow$\\[3pt]
Exact rational add / subtract / multiply\\[3pt]
$\downarrow$\\[3pt]
One output rounding and saturation\\[3pt]
$\downarrow$\\[3pt]
Result bits and diagnostic metadata
\end{minipage}}
\caption{Scalar arithmetic path implemented in phase I. The blocks describe
functional responsibilities, not hardware stages or cycle latency.}
\label{fig:architecture}
\end{figure}

The \code{TakumFormat} interface provides field extraction, encoding, decoding,
negation, and conversion between widths. The \code{TakumALU} interface applies
addition, subtraction, or multiplication to decoded operands and rounds the
exact result once. Successive operations each round their own output; no fused
operation or enlarged dot-product accumulator is implemented in this milestone.

Status metadata records inexact results, saturation below or above the
representable magnitude range, and $\NaR$. These are program diagnostics and
are not defined as RISC-V architectural CSRs. No pipeline, memory hierarchy,
vector register file, or instruction decoder is modeled yet.

\subsection{Verification strategy}
The first layer uses known bit patterns, round-trip checks, special-value
propagation, arithmetic identities, regime boundaries, width conversion, and CLI
tests. Every T8 and T16 encoding is decoded and re-encoded. T32 includes 10,000
seeded round-trip samples and additional boundary patterns.

The second layer supplies an independent exhaustive T8 oracle. It decodes
fields using a string-based implementation of the defining equations, enumerates
T9 codes at the rounding boundaries, and searches that ordered table. This
algorithm differs from constructing an encoding from an input exponent. It
checks all $256^2$ operand pairs for each of the three operations, including
$\NaR$, using exact rational arithmetic throughout.

The third layer compares with the pinned C reference. All T8/T16 patterns and
20,101 T32 patterns are checked for decoding. Encoding cases include decoded
values, 20,000 seeded binary floating-point inputs per width, and nonfinite or
zero inputs. Arithmetic is exhaustive in T8 and uses 20,000 random pairs plus
64 special/boundary pairs per operation in T16 and T32. Repeated inputs count
as comparisons, not distinct test values. The seed is 20260928.

\subsection{Reference portability correction}
The unmodified reference failed a basic smoke case on the local ARM Mac:
\code{takum32\_from\_float64(1.5)} returned \code{0x40000000}, encoding 1,
instead of \code{0x44000000}. In that environment, \code{long double} has the
precision of \code{double}. The reference's extended-float fraction extraction
increments its exponent after integer normalization, shifting the fraction
incorrectly. The routine also makes assumptions about the storage layout of
\code{long double}.

For validation, we replace that one extraction routine with a portable
calculation using \code{ldexpl}, \code{floorl}, and an explicit nearest/even tie
check. T32 requires at most 27 output fraction bits. The patch changes neither
the reference's characteristic decoding nor its three arithmetic operators.
The original checkout stays clean and a separate copy receives the patch.
The repository includes the patch, its hash in the validation report, the
upstream ISC license, and the original failed run. Accordingly, our reported
agreement is with the \emph{corrected} reference, not the unmodified library.

\section{Results and Discussion}
The local execution on an ARM Mac with Python 3.14.5 passed all 20 offline tests,
with no failures, errors, or skipped tests. These tests include 65,792 exhaustive
T8/T16 round-trip patterns and 196,608 T8 operations checked against the
independent exact oracle. Table~\ref{tab:validation} details the differential
coverage; all 548,793 comparisons passed with the portability correction.
The original reference produced 156,294 divergences on the same local test set,
which are retained as evidence of the portability problem.

\begin{table}[t]
\caption{Differential comparisons with the corrected reference}
\label{tab:validation}
\centering
\begin{tabular}{lrrr}
\toprule
Check & T8 & T16 & T32\\
\midrule
Decode & 256 & 65,536 & 20,101\\
Encode & 20,261 & 85,541 & 40,106\\
Add & 65,536 & 20,064 & 20,064\\
Subtract & 65,536 & 20,064 & 20,064\\
Multiply & 65,536 & 20,064 & 20,064\\
\midrule
Mismatches & 0 & 0 & 0\\
\bottomrule
\end{tabular}
\end{table}

A simple demonstration illustrates input quantization. The expression
$1.5+2.25$ produces $3.5$ in T8 because $2.25$ rounds to $2$ before addition.
In T16, both operands are exact and the result is $3.75$. Multiplying the
largest T8 value by itself saturates to that same largest value. Multiplying
the smallest positive T8 value by itself preserves a nonzero result through
lower saturation. These examples distinguish format behavior from errors in
the arithmetic implementation.

The exhaustive T8 oracle gives coverage over the complete operand domain for
the chosen operations and rounding policy. T16/T32 arithmetic coverage remains
sampled: agreement on these cases is not proof for all possible operands.
Additionally, the C reference computes through host floating point, while the
Python module computes exact rational intermediates. Carefully chosen inputs
outside the sample can therefore expose intermediate-rounding differences.
Rational boundary tests independently check the selected exact semantics.

No comparative SpMV results are reported at this stage. In particular, the
validation counts establish consistency of arithmetic behavior, not accuracy
advantages over FP8/FP16/FP32. Runtime measurements in the JSON records describe
the validation program and must not be interpreted as ALU latency or speedup.

\section{Conclusion and Next Steps}
The first milestone supplies a standalone linear Takum codec and scalar ALU for
T8, T16, and T32, with explicit special-value, saturation, and rounding semantics.
The combination of an independent exhaustive T8 oracle and a documented
external reference comparison makes the implementation reviewable and provides
a basis for subsequent integration.

The next milestone will specify vector registers, active vector length,
element width, masking, and the proposed arithmetic, dot-product, and narrowing
instructions. Later experiments will add sparse matrix-vector multiplication,
explicit FP8 baselines, a validated quadruple-precision reference, error
distributions, and instruction/CSR analysis. Those results will support the
final article and the 15-minute presentation.

\section*{Reproducibility}
Run \code{make verify} from the repository root. Offline tests require only
Python; differential tests additionally require Git, a C99 compiler, and network
access for the initial reference checkout. \code{make paper} compiles this
IEEE-format source using Tectonic. Machine-readable reports, the complete T8
value table, and CLI demonstration output are stored under \code{results/}.
The repository is available to the project collaborators at
\url{https://github.com/pedrohcali/aoc-takum-rvv}.

\begin{thebibliography}{3}
\bibitem{takum}
L. Hunhold, ``Beating Posits at Their Own Game: Takum Arithmetic,'' in
\emph{Next Generation Arithmetic}, LNCS, vol. 14666. Springer, 2024,
pp. 1 to 51, doi: 10.1007/978-3-031-72709-2\_1.
[Online]. Available: \url{https://arxiv.org/abs/2404.18603}
\bibitem{simd}
L. Hunhold, ``Streamlining SIMD ISA Extensions with Takum Arithmetic:
A Case Study on Intel AVX10.2,'' in \emph{2025 14th International Conference on
Modern Circuits and Systems Technologies (MOCAST)}, 2025, pp. 1 to 6.
[Online]. Available: \url{https://arxiv.org/abs/2503.14067}
\bibitem{libtakum}
L. Hunhold, ``libtakum: Takum arithmetic C99 reference implementation,''
version 2.0.0, revision \code{776ac0c} (full hash in the validation report).
[Online]. Available: \url{https://github.com/takum-arithmetic/libtakum}
\end{thebibliography}
\end{document}
